\chapter{Extended Background}
\label{chapter:Background}

This chapter provides the theoretical foundation for understanding the construction and evaluation of the trading strategies examined in this thesis. We begin by introducing technical trading indicators commonly used to extract information about trends, momentum, and market participation from historical price and volume data. These indicators form the basis for rule-based trading strategies, which are subsequently discussed as a transparent and structured approach to signal generation. Building on this foundation, we introduce machine learning–based signal filtering as a method for refining rule-generated signals while maintaining interpretability and controlling model complexity. Finally, the chapter presents key concepts related to backtesting and strategy evaluation, which establish the methodological framework for assessing performance, robustness, and generalization in the empirical analysis presented later in the thesis.

\section{Trading Indicators}

\subsection{Overview}

Technical indicators is a central component of rule-based and
algorithmic trading strategies, where they are used to extract information
from historical price and volume data in order to generate trading signals. \citep{curcio1997technical}.

Research showed that the performance of individual indicators is
often unstable across different market conditions, and lead the development
of more structured trading systems. As a result, later studies increasingly examined combinations of indicators in
order to capture multiple dimensions of market behaviour within a single
strategy
\citep{lento2007profitability}.

More recent empirical analyses compare a wide range of technical indicators,
including, for example EMA, RSI, and MACD
\citep{keshavarz2022trading}.
These indicators are commonly incorporated into rule-based trading systems,
where predefined rules translate indicator signals into buy and sell decisions
\citep{salkar2021algorithmic}.

No single indicator provide consistent results which lead to research combining multiple indicators to improve the accuracy of trading strategies
\citep{saud2024technical}.

In this study we will combine multiple technical indicators and use a multi-indicator approach as the foundational signal
generation layer of a rule-based trading strategy applied to the gold market.
Each selected indicator contributes complementary information related to price
trends, momentum dynamics, or market participation. The technical indicators that we will use are the EMA, RSI, MACD, DMI, KST and MFI, explained in more detail below.

%WRITE BELOW THIS COMMENT


\subsection{Exponential Moving Average}

The Exponential Moving Average (EMA) is a moving average of a series of historical price data over a specific time period. The EMA reacts faster to price movements as it gives higher weight to recent price action \citep{11203332}. 

The exponential moving average (EMA) at time \( t \) is defined as
\[
\mathrm{EMA}_t 
= \frac{2}{N+1} P_t 
+ \left(1 - \frac{2}{N+1}\right)\mathrm{EMA}_{t-1},
\]
where \( P_t \) denotes the asset price at time \( t \), and 
N represents the chosen window length.


In this thesis, the EMA will be used to construct a trend signal for the gold market.

\subsection{Relative Strength Index}

The Relative Strength Index (RSI) is a momentum-based technical indicator
(Salkar et al. 2021). The 14-period Relative Strength Index (RSI) is defined as
\[
RSI_t = 100 - \frac{100}{1 + RS_t},
\]
where
\[
RS_t = \frac{\overline{G}_t}{\overline{L}_t}.
\]
Here, $\overline{G}_t$ and $\overline{L}_t$ denote the exponentially smoothed
average gains and average losses over a 14-period lookback window, respectively.
Following Wilder’s original formulation, these quantities are updated as
\[
\overline{G}_t = \frac{13\,\overline{G}_{t-1} + G_t}{14}, \quad
\overline{L}_t = \frac{13\,\overline{L}_{t-1} + L_t}{14},
\]
where $G_t$ and $L_t$ denote the current period’s gain and loss
(Salkar et al. 2021).

Several recent studies indicate that RSI remains a relevant component in modern trading systems, particularly when combined with other indicators in multi-indicator strategies \citep{10795119}. 

RSI is used in this thesis to construct a rule-based signal that reflects momentum dynamics in the gold market.

\subsection{Moving Average Convergence Divergence}
The Moving Average Convergence Divergence (MACD) is a technical indicator
designed to capture changes in price momentum by measuring the relationship
between short-term and long-term exponential moving averages. Trading signals
are commonly derived from crossover events, where the MACD line crosses its
associated signal line, indicating potential changes in market direction
(Salkar et al. 2021).

The MACD at time $t$ is defined as
\[
MACD_t = EMA^{(\text{short})}_t - EMA^{(\text{long})}_t,
\]
where $EMA^{(\text{short})}_t$ and $EMA^{(\text{long})}_t$ denote exponential
moving averages computed using different window lengths in order to represent
distinct time horizons (Salkar et al. 2021). The signal line is commonly defined
as an exponential moving average of the MACD series.

Within the scope of this study, MACD serves as a momentum-oriented component
that complements trend- and volume-based indicators in a rule-based trading
framework.

\subsection{Directional Movement Index}

The Directional Movement Index (DMI) is designed to measure directional price movement by separating
positive and negative price changes over a given time period.
The indicator consists of two components, the positive directional indicator
(+DI) and the negative directional indicator (-DI), which represents upward and
downward price movements, respectively \citep{saud2024technical}.

The positive and negative directional indicators are defined as
\[
+DI_t = 100 \times \frac{EMA(+DM_t)}{ATR_t}, \quad
-DI_t = 100 \times \frac{EMA(-DM_t)}{ATR_t},
\]
where $+DM_t$ and $-DM_t$ denote positive and negative directional movement at
time $t$, $EMA(\cdot)$ represents exponential smoothing, and $ATR_t$ denotes
the average true range \citep{saud2024technical}.

In this thesis, DMI is included as a trend-oriented indicator that contributes
directional information to the rule-based trading system applied to the
gold market.

\subsection{Know Sure Thing}

The Know Sure Thing (KST) indicator is a momentum-based technical indicator
constructed by combining several rate-of-change (ROC) measures calculated
over different time horizons. Each ROC component is smoothed and weighted
before aggregation, allowing the indicator to capture longer-term price
momentum while reducing short-term fluctuations \citep{saud2024technical}.

The KST indicator at time $t$ is commonly defined as the weighted sum of a fixed
set of smoothed rate-of-change components,
\[
KST_t = \sum_{i=1}^{n} w_i \cdot \widetilde{ROC}_{i,t},
\]
where $\widetilde{ROC}_{i,t}$ denotes the smoothed rate of change over the
$i$-th time horizon and $w_i$ represents the corresponding weighting factor
\citep{saud2024technical}.

\subsection{Money Flow Index}

The Money Flow Index (MFI) is a momentum indicator that incorporates both
price and trading volume in its calculation. It is commonly interpreted as a
volume-weighted variant of the Relative Strength Index (RSI)
(Nugroho et al. 2014).

The money flow at time $t$ is defined as
\[
MF_t = TP_t \times V_t,
\]
where $TP_t = (H_t + L_t + C_t)/3$ denotes the typical price and $V_t$ denotes
the trading volume (Nugroho et al. 2014).

Over a selected lookback period, positive money flow is calculated when the
typical price increases compared to the previous period, while negative money
flow is calculated when the typical price decreases.

The money flow ratio is then computed as
\[
MFR_t = \frac{\text{Positive Money Flow}}{\text{Negative Money Flow}}.
\]

Finally, the Money Flow Index is defined as
\[
MFI_t = 100 - \frac{100}{1 + MFR_t}.
\]

The MFI ranges from 0 to 100. Values above 80 typically indicate
overbought conditions, while values below 20 indicate oversold conditions.

%START OF RULES SECTION

\section{Rule-Based Trading Strategies}

\subsection{Overview}

Rule-based trading strategies convert market data to trading actions such as \emph{buy}, \emph{sell}, or \emph{hold}. The background for evaluating such rules is the Efficient Market Hypothesis (EMH): if prices ``fully reflect'' available information, then systematic excess returns from publicly observable signals should be difficult to achieve after costs \citep{fama1970efficient}. 

Empirical research has investigated the performance of simple technical trading rules, such as moving-average–based strategies and filter rules, particularly in foreign exchange markets. Evidence from intraday foreign exchange data suggests that while apparent profitability may arise in certain market conditions, such profits are not robust on average once realistic transaction costs and trading constraints are incorporated, which is consistent with market efficiency \citep{curcio1997technical}. Similarly, genetic-programming-based searches for profitable rules on the gold market find that evolved rules can contain predictive content (e.g., identifying low-volatility, positive-return regimes), but do not deliver consistent excess returns over buy-and-hold after transaction costs in genuine out-of-sample tests \citep{allen1999using}. These results highlight a key practical point: even when rules extract some predictive structure, actual out-of-sample results and costs can erase statistical edges.

A complementary stream of work focuses on \emph{designing} and \emph{optimizing} rule sets. Evolutionary computation is frequently used to search the space of possible rule structures and parameterizations. Genetic programming has been proposed as a flexible framework for automatically generating trading rules, particularly when rules must adapt to individual assets rather than broad indices \citep{potvin2004generating}. Rule-based systems have also been combined with metaheuristic optimization to select and weight rule components, for example, using quantum-inspired tabu search together with sliding-window evaluation to mitigate overfitting \citep{chou2014rule}. Moreover, rule-based systems have been proposed that integrate fuzzy rules with evidential reasoning (Dempster--Shafer theory) to handle uncertainty and partially conflicting indicator evidence \citep{dymova2016forex}, while graph-based evolutionary methods such as Genetic Network Programming extract and accumulate buy/sell rules and then classify actions using rule pools and trend-dependent decision mechanisms \citep{mabu2013enhanced}.

In this thesis, rule-based strategies serve as a transparent \emph{signal generator}. We compute a set of indicator-based rules (EMA crossovers, RSI thresholds, and additional indicators such as MACD/DMI/KST as introduced in the previous section) and translate them into standardized trading signals for the gold market. These signals provide (i) a baseline aligned with classical technical analysis, and (ii) a structured feature set that can be combined into multi-indicator rules, following the general premise that agreement across signals may be more informative than any single rule \citep{lento2007profitability}. In subsequent chapters, we will treat these rule outputs as inputs to a machine-learning filtering stage, with the objective of down-weighting or rejecting low-quality signals while controlling for overfitting through out-of-sample evaluation.


\subsection{Signal Generation}

Trading rules are the basic components of the rule-based strategy. A trading rule is a predefined condition based on indicator values and/or the price series, and its purpose is to translate observed market information into a trading \emph{signal} \citep{ozturk2016heuristic}. In this thesis, signals take three forms: \emph{buy}, \emph{sell}, and \emph{hold}. A buy or sell signal is \emph{active} because it suggests taking action in the market, while a hold signal is \emph{passive} and means that no new action is taken.

A rule can be simple, for example, comparing an indicator to a fixed level, or more complex, for example, detecting specific crossover patterns. This generates a signal that can later be converted into a position \citep{ozturk2016heuristic}.

Rule-based signal generation is typically organized into families depending on how indicator information is used. The most common families are:

\begin{itemize}
    \item \textbf{Crossover rules}.These rules trigger signals when one time series crosses another (e.g., a fast-moving average crossing a slow-moving average), or when an indicator crosses a predefined threshold. In both cases, the crossing event is interpreted as a change in trend or momentum and is mapped into buy or sell signals \citep{ozturk2016heuristic}.
    \item \textbf{Divergence rules.} For momentum-type indicators, price and the indicator often move together. A divergence occurs when the price makes a new high or low, but the indicator does not confirm the move. Regular divergence is commonly interpreted as a potential trend reversal, while hidden divergence is commonly interpreted as trend continuation. These situations are converted into buy or sell signals depending on whether the divergence suggests upward or downward pressure \citep{ozturk2016heuristic}.
\end{itemize}

In this thesis, each indicator-based rule produces a standardized output signal (buy, sell, or hold) for each time step. These signals are used in two ways: (i) directly as a baseline rule-based strategy, and (ii) as structured inputs to the multi-indicator strategy and machine-learning filtering introduced later. This separation keeps the signal generation layer transparent and makes it easier to evaluate which parts of the overall pipeline add value in out-of-sample tests.



\subsection{Multi-Indicator Strategies}

A multi-indicator strategy is motivated by the idea that no single technical indicator is consistently reliable across markets, so combining indicators that capture different dimensions (e.g., trend, momentum, volatility, and volume) can provide a more holistic and interpretable basis for trading decisions \citep{10795119}. Further research shows that a multi-indicator approach outperforms a buy-and-hold strategy even when accounting for transaction costs \citep{lento2007profitability}.

In this thesis, we will therefore use a multi-indicator approach to increase our chances of creating a profitable strategy.


\section{Machine Learning-Based Signal Filtering}
\label{sec:ml_filtering}

\subsection{Overview of Machine Learning in Trading}

Machine learning methods have been applied in financial trading research as tools for pattern recognition and decision support based on historical market data \citep{dakalbab2024artificial}.

Empirical studies show that machine learning is often not used to replace trading rules entirely, but rather to enhance or refine rule-based systems by filtering or weighting generated signals \citep{hu2015stock}.
This hybrid approach aims to combine the interpretability of rule-based strategies with the adaptability of data-driven models \citep{dakalbab2024artificial}.

However, increased model flexibility also introduces a higher risk of overfitting \citep{dietterich1995overfitting}.
As a result, recent research emphasizes the use of relatively simple machine learning models when robustness and generalization are prioritized over predictive complexity \citep{bu2020research}.

In this study, machine learning is used as a post-processing step to evaluate signals generated by a predefined set of technical trading rules, with the model’s confidence score later used as the basis for performance analysis.

\subsection{Signal Filtering as a Classification Problem}
\label{subsec:signal_filtering_classification}

Signal filtering can be formulated as a supervised classification problem where candidate rule-generated trading points (signals) are evaluated using features derived from market data and technical analysis, and the target indicates whether the subsequent outcome is favorable according to a chosen criterion \citep{wu2006decisiontree, teixeira2010knn}.
In this setting, technical indicators and rule outputs can be treated as input attributes for a classifier that accepts or rejects candidate signals \citep{wu2006decisiontree}.

Several studies also model trading decisions as discrete classes (e.g., buy/sell/hold) to simplify learning and evaluation \citep{alnasseri2015stocktwits}.
This formulation allows models to learn which combinations of observed attributes are associated with desirable decision outcomes \citep{alnasseri2015stocktwits, wu2006decisiontree}.

Framing signal evaluation as classification supports transparent performance analysis using standard classification error measures and validation procedures. Such an approach is particularly suitable when the goal is to screen or filter candidate trading points produced by a rule, rather than to forecast an exact future price \citep{wu2006decisiontree, alnasseri2015stocktwits}.


\subsection{Decision Trees}
\label{subsec:decision_trees}

Building on the formulation of signal filtering as a classification task with discrete outcomes, the focus now turns to models that map indicator-based inputs to explicit decision rules. 

Decision trees have been applied in financial trading contexts to support
buy and sell decisions based on technical indicators and historical market
data. In this setting, trading decisions are formulated as a supervised learning
problem, where indicator-based observations are used to classify trading
signals into discrete outcomes. The tree-based structure enables trading decisions to be expressed as a
set of explicit decision rules, allowing inspection of which indicators
and threshold values influence the final decision \citep{nugroho2014decision}.

In financial applications, decision trees have been used to screen candidate trading points generated by technical rules such as filter rules \citep{wu2006decisiontree}.
Their structure enables inspection of which variables and thresholds drive accept/reject decisions \citep{wu2006decisiontree, alnasseri2015stocktwits}.

Decision trees can overfit without constraints, so depth control and validation are important in practice \citep{dietterich1995overfitting}.
This makes them suitable as filtering mechanisms in hybrid trading systems where transparency and controlled complexity are priorities \citep{wu2006decisiontree}.

In this thesis, a decision tree is used to evaluate rule-based trading signals on
the basis of indicator-derived features, with the resulting confidence scores used
to analyze how signal quality relates to subsequent trading performance.


\subsection{Overfitting and Generalization}
\label{subsec:overfitting_generalization}

Overfitting occurs when a model adapts to noise in training data rather than capturing structure that generalizes, leading to degraded out-of-sample performance \citep{dietterich1995overfitting, hawkins2004overfitting}.
For market data, changing market dynamics can further reduce a model's stability over time \citep{rechenthin2014ml}.

To mitigate overfitting and selection effects in strategy research, strict out-of-sample evaluation is required, and researchers highlight the risk of backtest overfitting when many strategies or variants are tried \citep{bailey2017pbo}.
In this study, these considerations motivate the use of a simple decision tree and strict separation between training and evaluation data.



%START OF BACKTESTING AND STRATEGY EVALUATION

\section{Backtesting and Strategy Evaluation}
\label{sec:backtesting}

\subsection{Overview of Backtesting}

Backtesting evaluates a trading strategy by simulating its decisions on historical data.
It is useful for comparing strategies under the same market conditions, but it can also
produce overly optimistic results if the research process is not controlled.

A key problem is \emph{data snooping}: trying many rules, parameter settings, or models
variants on the same dataset and then focusing on the best-looking outcome.
Even if markets were fully efficient, some strategy would appear to ``work'' by chance
when the search is large enough. White \citep{white2000reality} formalizes this issue
and proposes the \emph{Reality Check}, a bootstrap-based test that evaluates whether the
\emph{best} strategy found in a search truly beats a benchmark \emph{after accounting for
the search itself}. Without this correction, ``naive'' significance measures can be
misleading because they treat the chosen strategy as if it were the only one tested
\citep{white2000reality}.

Sullivan et al.\ \citep{sullivan1999data} apply this idea to technical trading rules by
testing a large rule universe and showing that adjusting for snooping is crucial when
interpreting apparent profitability. More generally, the message is strong in-sample
backtest results are not sufficient evidence of a real edge if they come from extensive
trial-and-error on the same data \citep{sullivan1999data, white2000reality}.

A second, related risk is \emph{backtest overfitting}: even with a train/test split, if
the strategy is selected because it performed best in-sample, the selected ``winner''
often disappoints out-of-sample. Bailey et al.\ \citep{bailey2017probability} propose the
\emph{probability of backtest overfitting} (PBO), estimated using repeated, structured
splits (combinatorially symmetric cross-validation), to quantify how likely the research
process is to pick a false positive.

In this thesis, backtesting is used to compare (i) single-indicator rule strategies,
(ii) multi-indicator rule strategies, and (iii) rule-based signals filtered by a machine
learning model. We use time-ordered evaluation, include transaction costs, and avoid any
tuning on the final test period. Results are interpreted with data-snooping and selection
risk in mind \citep{white2000reality, sullivan1999data, bailey2017probability}.


\subsection{Backtesting Framework}

In empirical finance, backtesting is used to evaluate trading strategies by
simulating their performance on historical data. A backtesting framework
defines how trading rules are assessed and how their performance is measured
over a given sample period \citep{white2000reality}.

A central challenge in backtesting is that the evaluation procedure itself can
influence the observed performance. White (2000) emphasizes that when many
strategies or rule variations are evaluated on the same dataset, conventional
performance measures may overstate statistical significance if the testing
process is not properly controlled \citep{white2000reality}.

Another important component of backtesting frameworks is the treatment of
transaction costs. Sullivan et al.~(1999) show that the apparent profitability of
technical trading strategies often diminishes or disappears once realistic
transaction costs are incorporated. This effect is particularly pronounced for
strategies with high trading frequency, where turnover-related costs can erode
returns \citep{sullivan1999data}.

More recent studies highlight that backtesting frameworks must also account for
the risk of overfitting arising from model selection and repeated testing.
Bailey et al.~(2017a) discuss how backtest results can be distorted when
strategies are chosen based on in-sample performance, motivating the use of
robust evaluation procedures that separate strategy development from final
performance assessment \citep{bailey2017probability}.

\subsection{Main Sources of Bias in Backtests}

Backtest bias refers to systematic effects that cause simulated strategy
performance to appear more favorable than what could realistically be achieved
in practice \citep{bailey2017probability}.

Look-ahead bias occurs when trading decisions are based on information that was
not available at the time the decision was made, such as indicator values that
implicitly incorporate future prices \citep{bailey2017probability}. The use of
such information leads to inflated performance estimates, since strategies are
allowed to benefit from foresight that would not be available in real trading
conditions \citep{bailey2017probability}.

Data snooping arises when a large number of trading rules are evaluated on the
same dataset and the best-performing rule is selected ex post
\citep{white2000reality}. White (2000) shows that under these circumstances,
standard significance measures can be misleading, as some strategies will
appear profitable purely due to chance rather than genuine predictive ability
\citep{white2000reality}.

Even when out-of-sample testing is employed, strategy selection is often guided
by in-sample performance, which introduces selection bias into the evaluation
process \citep{bailey2017probability}. Bailey et al.~(2017a) formalize this issue
as a backtest overfitting and show that strategies chosen based on in-sample
results frequently fail to generalize to unseen data \citep{bailey2017probability}.

\subsection{Validation and Evaluation Design}

Validation and evaluation design plays a central role in assessing whether
observed strategy performance reflects genuine predictive ability, or is the
result of overfitting \citep{bailey2017probability}. In particular, separating
strategy development from the final performance evaluation is necessary to avoid
biased estimates of out-of-sample returns \citep{bailey2017probability}.

A common approach discussed in the literature is to partition the available
data into distinct subsets used for model estimation, model selection, and
final evaluation \citep{bailey2017probability}. Bailey et al.~(2017a) emphasize
that such separation reduces the risk of unintentionally tailoring strategies
to the evaluation sample, which would otherwise inflate reported performance
\citep{bailey2017probability}.

Beyond a single static split, rolling or walk-forward evaluation procedures are
often considered in order to account for potential changes in market behavior
over time \citep{bailey2017probability}. By repeatedly re-estimating strategies
on past data and evaluating them on subsequent periods, walk-forward designs
provide multiple out-of-sample assessments and help reveal whether observed
performance is stable or the result of a one-off realization
\citep{bailey2017probability}.

\subsection{Performance Metrics}

The evaluation of trading strategies typically considers both profitability and
risk, since higher returns alone do not provide a complete picture of strategy
performance \citep{bailey2017probability}. Bailey et al.~(2017a) emphasize that
ignoring risk characteristics can lead to misleading conclusions about the
economic value of a strategy \citep{bailey2017probability}.

Return-based performance measures are commonly used to summarize the overall
profitability of a trading strategy over an evaluation period
\citep{bailey2017probability}. Such measures provide a basic assessment of how
strategy returns evolve over time, but do not capture variability or downside
risk on their own \citep{bailey2017probability}.

To account for risk, studies often complement return measures with indicators
of return variability and drawdowns \citep{bailey2017probability}. Bailey et
al.~(2017a) note that volatility-based measures and maximum drawdown are useful
for characterizing the stability of returns and the magnitude of potential
losses experienced by a strategy \citep{bailey2017probability}.

Risk-adjusted performance metrics combine information about returns and risk
into a single summary statistic \citep{bailey2017probability}. The Sharpe ratio
is one of the most commonly used risk-adjusted measures and expresses excess
return relative to return variability \citep{bailey2017probability}.

In addition to returns and risk, trading activity is relevant when evaluating
technical trading strategies, since transaction costs typically increase with
trading frequency \citep{sullivan1999data}. Sullivan et al.~(1999) show that
strategies generating frequent trades may appear profitable before costs but
often fail to maintain their performance once realistic transaction costs are
incorporated \citep{sullivan1999data}.

